\section{Group Actions and Permutation Representations}

\begin{theorem}
Let $G$ be any group and $A$ be a non-empty set. Then any homomorphism from $G$ to $\operatorname{Sym}(A)$, the symmetric group on $A$, defines an action of $G$ on $A$.

Conversely, every action of $G$ on $A$ induces a homomorphism from $G$ to $\operatorname{Sym}(A)$.
\end{theorem}

\begin{proof}
We establish the theorem in two parts:

\paragraph{Part 1: Homomorphism induces an action.}
Let $F : G \to \operatorname{Sym}(A)$ be a group homomorphism. For each $g \in G$, denote $F(g) = \sigma_g$, where $\sigma_g : A \to A$ is a bijective map. Since $F$ is a homomorphism, for any $g_1, g_2 \in G$:
\[
F(g_1 g_2) = F(g_1) \circ F(g_2) \implies \sigma_{g_1 g_2} = \sigma_{g_1} \circ \sigma_{g_2}.
\]
Define $* : G \times A \to A$ by:
\[
g * a = \sigma_g(a) \quad \forall g \in G, \; \forall a \in A.
\]
Then $*$ is a map. For any $g_1, g_2 \in G$ and $a \in A$:
\begin{align*}
g_1 * (g_2 * a) &= g_1 * (\sigma_{g_2}(a)) \\
&= \sigma_{g_1}(\sigma_{g_2}(a)) \\
&= (\sigma_{g_1} \circ \sigma_{g_2})(a) \\
&= \sigma_{g_1 g_2}(a) \\
&= (g_1 g_2) * a.
\end{align*}
Furthermore, since $F$ is a homomorphism, $F(e) = \sigma_e = I_A$, where $I_A$ is the identity permutation on $A$. Thus:
\[
e * a = \sigma_e(a) = I_A(a) = a.
\]
Therefore, $*$ is a group action of $G$ on $A$.

\paragraph{Part 2: Action induces a homomorphism.}
Conversely, suppose $* : G \times A \to A$ is a group action. Define $F : G \to \operatorname{Sym}(A)$ by setting $F(g) = \sigma_g$ for each $g \in G$, where $\sigma_g : A \to A$ is defined as:
\[
\sigma_g(a) = g * a \quad \forall a \in A.
\]
First, we show that $\sigma_g$ is a bijective map from $A$ to $A$:
\begin{itemize}
    \item \textbf{One-to-one (Injective):} Let $a_1, a_2 \in A$ such that $\sigma_g(a_1) = \sigma_g(a_2)$. Then:
    \begin{align*}
    g * a_1 = g * a_2 &\implies g^{-1} * (g * a_1) = g^{-1} * (g * a_2) \\
    &\implies (g^{-1} g) * a_1 = (g^{-1} g) * a_2 \quad (\text{as } * \text{ is an action}) \\
    &\implies e * a_1 = e * a_2 \\
    &\implies a_1 = a_2.
    \end{align*}
    Hence, $\sigma_g$ is one-to-one.
    
    \item \textbf{Onto (Surjective):} Let $a \in A$. Then $g^{-1} * a \in A$, and:
    \[
    \sigma_g(g^{-1} * a) = g * (g^{-1} * a) = (g g^{-1}) * a = e * a = a.
    \]
    Hence, $\sigma_g$ is onto.
\end{itemize}
Thus $\sigma_g$ is bijective, meaning $\sigma_g \in \operatorname{Sym}(A)$, and $F$ is a well-defined map.

Next, we show that $F$ is a group homomorphism. For any $g_1, g_2 \in G$ and $a \in A$:
\begin{align*}
\sigma_{g_1 g_2}(a) &= (g_1 g_2) * a \\
&= g_1 * (g_2 * a) \\
&= g_1 * \sigma_{g_2}(a) \\
&= \sigma_{g_1}(\sigma_{g_2}(a)) \\
&= (\sigma_{g_1} \circ \sigma_{g_2})(a).
\end{align*}
Therefore, $\sigma_{g_1 g_2} = \sigma_{g_1} \circ \sigma_{g_2}$, which implies:
\[
F(g_1 g_2) = F(g_1) \circ F(g_2).
\]
Hence, $F$ is a group homomorphism.
\end{proof}

\begin{remark}
For any group action of a group $G$ on a set $A$, the induced homomorphism
\[
F : G \longrightarrow \operatorname{Sym}(A)
\]
is called the \textbf{associated permutation representation} of the given action.
\end{remark}

\begin{corollary}[Cayley's Theorem]
Every group is isomorphic to a subgroup of its permutation group (symmetric group).
\end{corollary}

\begin{proof}
Let $G$ be any group. We know that $G$ acts on itself ($A = G$) under the left regular action:
\[
* : G \times A \longrightarrow A, \quad g * a = ga \quad \forall g \in G, \; \forall a \in A = G.
\]
By the previous theorem, the corresponding permutation representation is the homomorphism:
\[
F : G \longrightarrow \operatorname{Sym}(G), \quad F(g) = \sigma_g, \quad \text{where } \sigma_g(a) = ga.
\]
We show that $F$ is one-to-one:
\begin{align*}
F(g_1) = F(g_2) &\implies \sigma_{g_1} = \sigma_{g_2} \\
&\implies \sigma_{g_1}(a) = \sigma_{g_2}(a) \quad \forall a \in G \\
&\implies g_1 a = g_2 a \quad \forall a \in G.
\end{align*}
Choosing $a = e$:
\[
g_1 e = g_2 e \implies g_1 = g_2.
\]
Therefore, $F$ is one-to-one, which means:
\[
\ker F = \{e\}.
\]
By the Fundamental Theorem of Group Homomorphisms (First Isomorphism Theorem):
\[
\frac{G}{\ker F} \cong F(G) \implies \frac{G}{\{e\}} \cong F(G) \implies G \cong F(G).
\]
Since $F(G) \le \operatorname{Sym}(G)$, every group $G$ is isomorphic to a subgroup of $\operatorname{Sym}(G)$.
\end{proof}
